\documentclass{article}
\usepackage[T1]{fontenc}
\usepackage[utf8]{inputenc}
\usepackage[italian]{babel}
\usepackage[accsupp]{axessibility}
\usepackage{tgheros}
\renewcommand{\familydefault}{\sfdefault}
\usepackage[defaultmathsizes]{mathastext}
\usepackage{microtype}
\DisableLigatures{encoding = *, family = *}
\usepackage{ragged2e}
\RaggedRight
\hyphenpenalty=10000
\exhyphenpenalty=10000
\usepackage{setspace}
\setstretch{1.5}
\usepackage{titlesec}
\titleformat{\section}
    {\fontsize{18}{27}\selectfont\bfseries}
    {}{0em}{}

\newcommand{\sqrtIT}[1]{%
  \BeginAccSupp{method=escape,ActualText={radice quadrata di #1}}%
  \sqrt{#1}%
  \EndAccSupp{}%
}

% Macro per inline math con testo alternativo esplicito.
% ActualText={...} con graffe evita che le virgole vengano
% interpretate dal parser keyval di accsupp.
\newcommand{\inlineformula}[2]{%
  \BeginAccSupp{method=escape,ActualText={#1}}%
  $#2$%
  \EndAccSupp{}%
}

\AtBeginDocument{\pdfcatalog{/Lang (it-IT)}}
\pagestyle{empty}

\begin{document}
\fontsize{18}{27}\selectfont

\section*{Domanda 1}

Luigi ha 2 figli di 15 e 11 anni.
Fra 18 anni la sua età sarà uguale alla somma delle età
che avranno i suoi figli.
Quanti anni ha oggi Luigi?

\begin{enumerate}
    \item[a)] 30
    \item[b)] Non si può dire
    \item[c)] 52
    \item[d)] 44
\end{enumerate}

\textbf{Risposta esatta: (d)}

\newpage

\section*{Domanda 2}

Il resto della divisione del polinomio
\inlineformula{x alla quinta meno 3 x alla quarta più 3}{x^5 - 3x^4 + 3}
per \inlineformula{x più 1}{x + 1} è:

\begin{enumerate}
    \item[a)] 1
    \item[b)] \inlineformula{meno 1}{-1}
    \item[c)] 3
    \item[d)] 0
\end{enumerate}

\textbf{Risposta esatta: (b)}

\newpage

\section*{Domanda 3}

Sia \inlineformula{T}{T} un triangolo rettangolo isoscele.
Allora la somma dei coseni degli angoli interni di \inlineformula{T}{T}
è uguale a:

\begin{enumerate}
    \item[a)] 2
    \item[b)] 1
    \item[c)] $\sqrtIT{3}$
    \item[d)] $\sqrtIT{4}$
\end{enumerate}

\textbf{Risposta esatta: (d)}

\newpage

\section*{Domanda 4}

Il 30\% degli studenti iscritti a un corso universitario
ha superato l'esame relativo al primo appello.
Se, dei restanti studenti iscritti,
il 10\% supera l'esame al secondo appello,
gli studenti che devono ancora superare l'esame
dopo i primi due appelli saranno:

\begin{enumerate}
    \item[a)] Il 37\% del numero totale di studenti iscritti al corso
    \item[b)] Il 60\% del numero totale di studenti iscritti al corso
    \item[c)] Il 63\% del numero totale di studenti iscritti al corso
    \item[d)] Il 40\% del numero totale di studenti iscritti al corso
\end{enumerate}

\textbf{Risposta esatta: (c)}

\newpage

\section*{Domanda 5}

Trova il risultato della seguente disequazione:
\inlineformula{2 diviso x maggiore di 1}{\dfrac{2}{x} > 1}

\begin{enumerate}
    \item[a)] \inlineformula{x minore di 2}{x < 2}
    \item[b)] \inlineformula{0 minore di x minore di 2}{0 < x < 2}
    \item[c)] \inlineformula{x maggiore di 2}{x > 2}
    \item[d)] \inlineformula{0 minore di x}{0 < x} e \inlineformula{x maggiore di 2}{x > 2}
\end{enumerate}

\textbf{Risposta esatta: (b)}

\end{document}